\section*{Hoofdstuk \ref{ch:b1:lewis-resonance-vsepr} --- Lewisstructuren, resonantie en VSEPR}

\begin{solution}{exo:b1:lewis-resonance-vsepr:1}
\ce{H2O}: twee \ce{O-H}-bindingen, twee \omterm{def:b1:lewis-resonance-vsepr:lewis}{vrije paren} op O. \ce{NH3}: drie
\ce{N-H}-bindingen, één \omterm{def:b1:lewis-resonance-vsepr:lewis}{vrij paar} op N. \ce{CO2}: \ce{O=C=O}, twee \omterm{def:b1:lewis-resonance-vsepr:lewis}{vrije
paren} op elk O-atoom. \ce{HCN}: \ce{H-C#N}, één \omterm{def:b1:lewis-resonance-vsepr:lewis}{vrij paar} op N.
\ce{CH2O}: koolstof gebonden aan twee H en dubbel gebonden aan O, twee
\omterm{def:b1:lewis-resonance-vsepr:lewis}{vrije paren} op O. In deze vijf zijn alle \omterm{def:b1:lewis-resonance-vsepr:formal-charge}{formele ladingen} nul.
\ce{NH4+}: vier \ce{N-H}-bindingen, geen \omterm{def:b1:lewis-resonance-vsepr:lewis}{vrij paar}; stikstof
$5 - 0 - 4 = +1$, de lading van het ion.
\end{solution}

\begin{solution}{exo:b1:lewis-resonance-vsepr:2}
\ce{H2S}: \ce{AX2E2}, gebogen, onder $109.5^\circ$ (gemeten $92^\circ$).
\ce{PCl3}: \ce{AX3E}, trigonaal piramidaal, ongeveer $100^\circ$.
\ce{BF3}: \ce{AX3}, trigonaal vlak, $120^\circ$. \ce{SiH4}: \ce{AX4},
tetraëdrisch, $109.5^\circ$. \ce{CO3^{2-}}: \ce{AX3} (de dubbele binding
telt één keer), trigonaal vlak, $120^\circ$.
% ledger: geo:H2S.aHSH, geo:PCl3.aClPCl
\end{solution}

\begin{solution}{exo:b1:lewis-resonance-vsepr:3}
Apolair: \ce{CCl4} (tetraëdrisch \ce{AX4}), \ce{BF3} (\ce{AX3}),
\ce{CO2} (lineair \ce{AX2}): de \omterm{def:b1:lewis-resonance-vsepr:dipole}{bindingsdipolen} heffen elkaar op. Polair:
\ce{CH2Cl2} (twee verschillende soorten bindingen op een tetraëder),
\ce{NH3} (piramidaal), \ce{SO2} (gebogen).
\end{solution}

\begin{solution}{exo:b1:lewis-resonance-vsepr:4}
\ce{H+} is het \omterm{def:b1:lewis-resonance-vsepr:lewis-acid}{lewiszuur} en \ce{H2O} de base: de nieuwe \ce{O-H}-binding
van \ce{H3O+} is bij haar vorming datief (daarna zijn de drie
\ce{O-H}-bindingen identiek). \ce{Cu^{2+}} is het zuur en elk \ce{NH3}
een base: de vier \ce{Cu-N}-bindingen zijn datief.
\end{solution}

\begin{solution}{exo:b1:lewis-resonance-vsepr:5}
\chemfig{CH_3-C(=[:60]O)-[:-60]\chemabove{O}{\scriptstyle\ominus}}
$\leftrightarrow$
\chemfig{CH_3-C(-[:60]\chemabove{O}{\scriptstyle\ominus})=[:-60]O}.
De twee structuren zijn gelijkwaardig: beide \ce{C-O}-bindingen hebben
orde 1.5 en dezelfde lengte, tussen de dubbele en de enkele binding in.
In methaanzuur verschillen de twee zuurstofatomen (één draagt de H), de
structuren zijn niet gelijkwaardig en de bindingen houden verschillende
lengten, 120.2 (\ce{C=O}) en \qty{134.3}{pm} (\ce{C-OH}).
\end{solution}

\begin{solution}{exo:b1:lewis-resonance-vsepr:6}
$N = 6 + 4 \times 6 + 2 = 32$ elektronen. Met vier enkele bindingen en
drie \omterm{def:b1:lewis-resonance-vsepr:lewis}{vrije paren} op elk zuurstofatoom: zwavel $6 - 0 - 4 = +2$, elk
zuurstofatoom $6 - 6 - 1 = -1$; totaal $+2 - 4 = -2$. Vorm \ce{AX4},
tetraëdrisch. Met twee \ce{S=O}-bindingen geeft de keuze van de twee van
de vier zuurstofatomen die dubbel gebonden zijn $\binom{4}{2} = 6$
structuren.
\end{solution}

\begin{solution}{exo:b1:lewis-resonance-vsepr:7}
\ce{SF4}: \ce{AX4E}, wipplank. \ce{ClF3}: \ce{AX3E2}, T-vormig.
\ce{XeF2}: \ce{AX2E3}, lineair. \ce{XeF4}: \ce{AX4E2}, vierkant vlak. In
een trigonale bipiramide heeft een axiale positie drie buren op
$90^\circ$, een equatoriale positie maar twee: de omvangrijke \omterm{def:b1:lewis-resonance-vsepr:lewis}{vrije paren}
gaan daar zitten waar ze de minste afstotingen op $90^\circ$ ondervinden.
\end{solution}

\begin{solution}{exo:b1:lewis-resonance-vsepr:8}
Zwavel en fosfor zijn groter en minder elektronegatief dan zuurstof en
stikstof: de \omterm{def:b1:lewis-resonance-vsepr:lewis}{bindende paren} liggen verder van het centrale atoom en worden
naar waterstof getrokken, zodat ze elkaar minder afstoten, en de \omterm{def:b1:lewis-resonance-vsepr:lewis}{vrije
paren} drukken de bindingen naar $90^\circ$.
\end{solution}

\begin{solution}{exo:b1:lewis-resonance-vsepr:9}
\ce{CH3-C#N}: het koolstofatoom van \ce{CH3} is sp$^3$, het
nitrilkoolstofatoom sp, het stikstofatoom sp; 5 $\sigma$-bindingen (drie
\ce{C-H}, \ce{C-C}, één in \ce{C#N}) en 2 $\pi$-bindingen.
\ce{CH2=CH-CH3}: de twee alkeenkoolstofatomen zijn sp$^2$, het
methylkoolstofatoom sp$^3$; 8 $\sigma$-bindingen (zes \ce{C-H}, twee
\ce{C-C}) en 1 $\pi$-binding.
\end{solution}

\begin{solution}{exo:b1:lewis-resonance-vsepr:10}
Beide moleculen zijn piramidaal met een \omterm{def:b1:lewis-resonance-vsepr:lewis}{vrij paar} op stikstof. Het \omterm{def:b1:lewis-resonance-vsepr:lewis}{vrije
paar} levert een bijdrage die (van negatief naar positief) van het \omterm{def:b1:lewis-resonance-vsepr:lewis}{vrije
paar} naar de kern wijst. In \ce{NH3} is stikstof het negatieve uiteinde
van elke binding: de \omterm{def:b1:lewis-resonance-vsepr:dipole}{bindingsdipolen} wijzen van N naar de
waterstofatomen, in dezelfde richting als de bijdrage van het \omterm{def:b1:lewis-resonance-vsepr:lewis}{vrije
paar}, en ze tellen op. In \ce{NF3} is fluor het negatieve uiteinde: de
\omterm{def:b1:lewis-resonance-vsepr:dipole}{bindingsdipolen} wijzen van de fluoratomen naar N, tegen de bijdrage van
het \omterm{def:b1:lewis-resonance-vsepr:lewis}{vrije paar} in, en de twee heffen elkaar bijna op.
\end{solution}

\begin{solution}{exo:b1:lewis-resonance-vsepr:11}
$\mu_b = \mu/(2\cos(\theta/2)) = 1.63/(2\cos 59.75^\circ) =
\qty{1.62}{\debye} = \qty{5.40e-30}{C.m}$. Dan is $\delta = \mu_b/(e\,d) =
5.40 \times 10^{-30}/(1.602 \times 10^{-19} \times 143.2 \times 10^{-12})
= 0.24$.
\end{solution}

\begin{solution}{exo:b1:lewis-resonance-vsepr:12}
\ce{N=N=O} met \omterm{def:b1:lewis-resonance-vsepr:formal-charge}{formele ladingen} $-1$, $+1$, $0$; \ce{N#N-O} met $0$,
$+1$, $-1$. De \ce{N-N}-binding (\qty{112.8}{pm}) ligt dicht bij de
drievoudige binding van \ce{N2} (\qty{109.8}{pm}): de structuur
\ce{N#N-O} weegt het zwaarst, en ze legt de negatieve lading op zuurstof,
het elektronegatiefste atoom. De \ce{N-O}-binding is echter korter dan
een enkele binding: de andere structuur draagt ook bij.
\end{solution}

\begin{solution}{pb:b1:lewis-resonance-vsepr:1}
\textbf{1.} \ce{N2} 10, \ce{NO} 11, \ce{NO2} 17, \ce{N2O} 16, \ce{O3} 18,
\ce{HNO3} 24.
\textbf{2.} \ce{N#N} met één \omterm{def:b1:lewis-resonance-vsepr:lewis}{vrij paar} op elk stikstofatoom.
\textbf{3.} \ce{N=O} met twee \omterm{def:b1:lewis-resonance-vsepr:lewis}{vrije paren} op zuurstof, en één \omterm{def:b1:lewis-resonance-vsepr:lewis}{vrij paar}
en één enkel elektron op stikstof. Met 11 elektronen, een oneven aantal,
heeft altijd één atoom een oneven telling: stikstof heeft 7 elektronen om
zich heen.
\textbf{4.} \ce{N=N=O}: $-1$, $+1$, $0$; \ce{N#N-O}: $0$, $+1$, $-1$.
\textbf{5.} \ce{O=O-O}: centraal zuurstofatoom $+1$ (één \omterm{def:b1:lewis-resonance-vsepr:lewis}{vrij paar}, drie
bindingen), het enkelvoudig gebonden eindstandige zuurstofatoom $-1$, het
andere 0.
\textbf{6.} Zoals in \cref{ex:b1:lewis-resonance-vsepr:hno3}: stikstof
$+1$, het zuurstofatoom met een enkele binding en zonder waterstof $-1$,
de andere 0.
\textbf{7.} \ce{NO} en \ce{NO2} (11 en 17 elektronen).
\textbf{8.} Twee gelijkwaardige structuren, met de dubbele binding aan
de ene of de andere kant: orde 1.5 voor elke binding.
\textbf{9.} Ja: \qty{127.8}{pm} ligt tussen de dubbele binding
(\qty{120.8}{pm}) en de enkele binding (\qty{147.5}{pm}) in.
\textbf{10.} De dubbele binding beurtelings op elk van de drie
zuurstofatomen; elke \ce{N-O}-binding is dubbel in één structuur op drie:
orde $4/3$.
\textbf{11.} De lange binding, \qty{140.6}{pm}, is \ce{N-OH}, een enkele
binding. De twee zuurstofatomen zonder waterstof delen één
$\pi$-binding via twee gelijkwaardige \omterm{def:b1:lewis-resonance-vsepr:resonance}{resonantiestructuren}: twee
bindingen van orde 1.5, 121.1 en \qty{119.9}{pm}.
\textbf{12.} Stikstof draagt een enkel elektron in plaats van een \omterm{def:b1:lewis-resonance-vsepr:lewis}{vrij
paar}: dat stoot de \omterm{def:b1:lewis-resonance-vsepr:lewis}{bindende paren} minder af dan een paar zou doen, en de
bindingen openen zich voorbij $120^\circ$.
\textbf{13.} \ce{N#N-O}, met de negatieve lading op zuurstof.
\textbf{14.} \ce{O3}: \ce{AX2E}, gebogen. \ce{N2O}: \ce{AX2} (centraal
N), lineair. \ce{NO3-}: \ce{AX3}, trigonaal vlak. Stikstof van
\ce{HNO3}: \ce{AX3}, trigonaal vlak.
\textbf{15.} Het \omterm{def:b1:lewis-resonance-vsepr:lewis}{vrije paar} van het centrale zuurstofatoom stoot de
\omterm{def:b1:lewis-resonance-vsepr:lewis}{bindende paren} sterker af dan die elkaar afstoten.
\textbf{16.} Ozon is gebogen en zijn centrale atoom verschilt van de
eindatomen (positieve \omterm{def:b1:lewis-resonance-vsepr:formal-charge}{formele lading} in het midden, negatieve lading
verdeeld over de uiteinden): een kleine dipool langs de bissectrice.
\ce{N2O} is lineair, maar zijn twee uiteinden zijn verschillende atomen:
een kleine dipool. \ce{N2} is symmetrisch: geen dipool.
\textbf{17.} \ce{CO2} is lineair (\ce{AX2}): zijn \omterm{def:b1:lewis-resonance-vsepr:dipole}{bindingsdipolen} heffen
elkaar op. \ce{SO2} is gebogen (\ce{AX2E}): ze tellen op.
\textbf{18.} \ce{AX3E}: het \omterm{def:b1:lewis-resonance-vsepr:lewis}{vrije paar} verkleint de hoeken tot onder
$109.5^\circ$; gemeten $106.7^\circ$.
\textbf{19.} \ce{AX2E2}: twee \omterm{def:b1:lewis-resonance-vsepr:lewis}{vrije paren} verkleinen ze nog meer;
$104.5^\circ$.
\textbf{20.} $\mu = 2\mu_{OH}\cos(\theta/2)$
(\cref{prop:b1:lewis-resonance-vsepr:dipole-sum}).
\textbf{21.} $\mu_{OH} = 1.857/(2\cos 52.24^\circ) = 1.857/1.2246 =
\qty{1.52}{\debye}$.
\textbf{22.} $1.516 \times 3.336 \times 10^{-30} = \qty{5.06e-30}{C.m}$.
\textbf{23.} $e\,d = 1.602 \times 10^{-19} \times 95.8 \times 10^{-12} =
\qty{1.535e-29}{C.m}$, dat is \qty{4.60}{\debye}.
\textbf{24.} $\delta = 5.06 \times 10^{-30}/1.535 \times 10^{-29} =
\textbf{0.33}$: de \ce{O-H}-binding van water draagt een derde van de
lading die een volledig ionaire binding zou dragen.
\end{solution}
